\documentclass[12pt,openany]{book}

% ============================================================
% Pequeno Livro de Linux para Estudantes de Estatística e
% Ciência de Dados
% Arquivo principal: main.tex
% ============================================================

% ----------------------------
% Idioma, codificação e fontes
% ----------------------------
\usepackage[utf8]{inputenc}
\usepackage[T1]{fontenc}
\usepackage[brazil]{babel}
\usepackage{lmodern}
\usepackage{microtype}

% ----------------------------
% Geometria e layout
% ----------------------------
\usepackage[a4paper,margin=2.8cm]{geometry}
\usepackage{setspace}
\onehalfspacing

% ----------------------------
% Matemática, tabelas e figuras
% ----------------------------
\usepackage{amsmath,amssymb}
\usepackage{graphicx}
\usepackage{booktabs}
\usepackage{longtable}
\usepackage{array}

% ----------------------------
% Cores, links e metadados
% ----------------------------
\usepackage{xcolor}
\definecolor{gulecdblue}{RGB}{0,85,150}

\usepackage[
    colorlinks=true,
    linkcolor=gulecdblue,
    citecolor=gulecdblue,
    urlcolor=gulecdblue
]{hyperref}

% ----------------------------
% Bibliografia
% ----------------------------
\usepackage[
    backend=biber,
    style=authoryear,
    sorting=nyt,
    refsection=chapter
]{biblatex}

% Adicione o seu arquivo bib aqui em uma nova linha definida por \addbibresource
\addbibresource{ref_marcosdecarvalho.bib}

% Categorias para leituras recomendadas por parte.
% As chaves devem existir nos arquivos bib definidos acima
\DeclareBibliographyCategory{recomendadosParteI}
\DeclareBibliographyCategory{recomendadosParteII}
\DeclareBibliographyCategory{recomendadosParteIII}
\DeclareBibliographyCategory{recomendadosParteIV}
\DeclareBibliographyCategory{recomendadosParteV}
\DeclareBibliographyCategory{recomendadosParteVI}
\DeclareBibliographyCategory{recomendadosParteVII}
\DeclareBibliographyCategory{recomendadosProximosPassos}

\addtocategory{recomendadosParteI}{lacroixLinuxMintEssentials2014,jrLinuxCommandLine2012}
\addtocategory{recomendadosParteII}{flyntLinuxShellScripting2017}
\addtocategory{recomendadosParteIII}{flyntLinuxShellScripting2017}
\addtocategory{recomendadosParteIV}{flyntLinuxShellScripting2017}
\addtocategory{recomendadosParteV}{flyntLinuxShellScripting2017}
\addtocategory{recomendadosParteVI}{flyntLinuxShellScripting2017}
\addtocategory{recomendadosParteVII}{flyntLinuxShellScripting2017}
\addtocategory{recomendadosProximosPassos}{flyntLinuxShellScripting2017}

% ----------------------------
% Glossário
% ----------------------------
\usepackage[toc,acronym]{glossaries}
\makeglossaries

\newglossaryentry{kernel}{
    name={kernel},
    description={Núcleo do sistema operacional responsável pela mediação entre hardware, processos, memória, dispositivos e chamadas de sistema}
}

\newglossaryentry{shell}{
    name={shell},
    description={Programa que interpreta comandos do usuário e permite interação textual com o sistema operacional}
}

\newglossaryentry{distro}{
    name={distribuição Linux},
    description={Sistema operacional completo construído a partir do kernel Linux, utilitários, bibliotecas, gerenciador de pacotes, ambiente gráfico e políticas de manutenção}
}

\newglossaryentry{pipe}{
    name={pipe},
    description={Mecanismo de composição em que a saída de um processo é conectada à entrada de outro processo}
}

\newglossaryentry{reprodutibilidade}{
    name={reprodutibilidade computacional},
    description={Capacidade de executar novamente um fluxo computacional e obter resultados equivalentes sob condições controladas e documentadas}
}

\newglossaryentry{hpc}{
    name={HPC},
    description={Computação de alto desempenho, geralmente realizada em clusters, servidores ou ambientes paralelos especializados}
}

\newglossaryentry{cli}{
    name={CLI},
    description={Interface de linha de comando, do inglês command-line interface}
}

\newglossaryentry{gui}{
    name={GUI},
    description={Interface gráfica de usuário, do inglês graphical user interface}
}

% ----------------------------
% Comandos auxiliares
% ----------------------------
\newcommand{\placeholder}[1]{
    \par\medskip
    \noindent\textit{#1}
    \par\medskip
}

\newcommand{\topicos}[1]{
    \begin{itemize}
        #1
    \end{itemize}
}

\newcommand{\bibliografiaParte}[1]{
    \cleardoublepage
    \section*{Bibliografia da Parte}
    \addcontentsline{toc}{section}{Bibliografia da Parte}
    \printbibliography[
        heading=none,
        notcategory=#1
    ]

    \section*{Leitura Recomendada}
    \addcontentsline{toc}{section}{Leitura Recomendada}
    \nocite{*}
    \printbibliography[
        heading=none,
        category=#1
    ]
}

% ============================================================
% Início do documento
% ============================================================

\begin{document}

% ============================================================
% Capa externa
% ============================================================

\begin{titlepage}
    \thispagestyle{empty}
    \centering

    \vspace*{3cm}

    {\Huge\bfseries Pequeno Livro de Linux\par}
    \vspace{0.5cm}
    {\Large\bfseries para Estudantes de Estatística e Ciência de Dados\par}

    \vspace{3cm}

    {\Large Grupo de Usuários Linux e Software Livre\par}
    {\Large do Curso de Estatística e Ciência de Dados da UFPR\par}

    \vfill

    \placeholder{Espaço reservado para a futura arte de capa, ilustração, gráfico ou composição visual do livro.}

    \vfill

    {\large Curitiba, Paraná\par}
    {\large \today\par}
\end{titlepage}

% ============================================================
% Capa interna
% ============================================================

\cleardoublepage
\pagenumbering{arabic}
\setcounter{page}{2}

\begin{titlepage}
    \thispagestyle{plain}
    \centering

    \vspace*{6cm}

    {\Huge\bfseries Pequeno Livro de Linux\par}
    \vspace{0.5cm}
    {\Large\bfseries para Estudantes de Estatística e Ciência de Dados\par}

    \vfill
\end{titlepage}

% ============================================================
% Autores
% ============================================================

\cleardoublepage
\chapter*{Autores}
\addcontentsline{toc}{chapter}{Autores}

Este livro é uma obra colaborativa produzida no contexto do Grupo de Usuários Linux e Software Livre do Curso de Estatística e Ciência de Dados da UFPR.

\section*{Lista de autores}

%%%%%%%%%%%%%%%%%%%%%%%%%%%%%%%%%%%%%%%%%%%%%%%%%%%%%%%%%%%%%%%%%%%%%%%%%%%%%%%%%%%%%%%%%%%%%%%%%%%%%%%%%%%%%%%%%%%%%%%%%%%
%%%%%%%%%%%%%%%%%%%%%%%%%%%%%%%%%%%%%%%%%%%%%%%%%%%%%%%%%%%%%%%%%%%%%%%%%%%%%%%%%%%%%%%%%%%%%%%%%%%%%%%%%%%%%%%%%%%%%%%%%%%

% Os colaboradores que realizarem contribuições substanciais ao conteúdo do livro devem adicionar seus nomes nesta seção.

%%%%%%%%%%%%%%%%%%%%%%%%%%%%%%%%%%%%%%%%%%%%%%%%%%%%%%%%%%%%%%%%%%%%%%%%%%%%%%%%%%%%%%%%%%%%%%%%%%%%%%%%%%%%%%%%%%%%%%%%%%%
%%%%%%%%%%%%%%%%%%%%%%%%%%%%%%%%%%%%%%%%%%%%%%%%%%%%%%%%%%%%%%%%%%%%%%%%%%%%%%%%%%%%%%%%%%%%%%%%%%%%%%%%%%%%%%%%%%%%%%%%%%%

\begin{itemize}
    \item Marcos de Carvalho
\end{itemize}

% ============================================================
% Agradecimentos
% ============================================================

\cleardoublepage
\chapter*{Agradecimentos}
\addcontentsline{toc}{chapter}{Agradecimentos}

Agradecemos aos estudantes, professores, técnicos, colaboradores e membros da comunidade de software livre que contribuírem para a construção deste material. Este livro nasce da convicção de que o conhecimento técnico deve ser compartilhado, criticado, melhorado e colocado a serviço da formação científica.

Agradecemos também ao Curso de Estatística e Ciência de Dados da Universidade Federal do Paraná, ao Departamento de Estatística e à comunidade acadêmica que apoia iniciativas de formação aberta e colaborativa.

% ============================================================
% Licença
% ============================================================

\cleardoublepage
\chapter*{Licença}
\addcontentsline{toc}{chapter}{Licença}

Este livro é distribuído sob os termos da licença \textbf{GNU General Public License v3.0}, também conhecida como \textbf{GPL-3.0}.

A versão juridicamente oficial da GPL-3.0 é o texto original em inglês publicado pela Free Software Foundation. O texto completo pode ser encontrado na raiz do repositório deste livro, no arquivo \texttt{LICENSE}. % e também abaixo.

% TODO -- \section*{Texto oficial da licença}

%\placeholder{Inserir aqui, se desejado, o texto oficial integral da GNU General Public License v3.0 em inglês por meio de um arquivo externo, por exemplo: \texttt{\textbackslash input\{licenca/gpl-3.0.tex\}}.}

\section*{Descrição da GPL-3.0 no contexto deste livro}

No contexto deste livro, a GPL-3.0 significa que qualquer pessoa pode ler, copiar, redistribuir, estudar, modificar e criar versões derivadas da obra, desde que as versões modificadas também sejam distribuídas sob a mesma licença. Isso preserva a liberdade do material ao longo do tempo e impede que uma versão derivada seja fechada ou redistribuída de forma incompatível com os princípios de acesso aberto e colaboração.

Como o livro é escrito em LaTeX, o código-fonte da obra corresponde aos arquivos \texttt{.tex}, aos arquivos bibliográficos \texttt{.bib}, às figuras, tabelas, scripts auxiliares e demais arquivos necessários para reconstruir a versão final do livro. Assim, qualquer redistribuição substancial da obra deve preservar o acesso a esses arquivos-fonte.

A licença também exige que os autores originais sejam creditados adequadamente. Portanto, contribuições relevantes devem ser registradas na lista de autores deste documento, bem como no histórico de contribuições do repositório.

% ============================================================
% Preâmbulo
% ============================================================

\cleardoublepage
\chapter*{Preâmbulo}
\addcontentsline{toc}{chapter}{Preâmbulo}

Este livro surgiu da necessidade de produzir uma introdução ao Linux que fosse relevante para estudantes de Estatística e Ciência de Dados. Muitos materiais introdutórios tratam o Linux como um conjunto isolado de comandos, sem mostrar sua importância em processos de análise de dados, pesquisa científica, automação, reprodutibilidade e computação moderna. Este projeto parte da premissa oposta: Linux é uma infraestrutura intelectual e técnica para pensar, organizar, executar e reproduzir trabalho computacional.

O público-alvo principal são estudantes de graduação em Estatística e Ciência de Dados, especialmente aqueles que estão iniciando seu contato com sistemas Unix-like, linha de comando, ambientes computacionais científicos, programação, análise de dados e ferramentas de reprodutibilidade. O texto também pode ser útil para estudantes de pós-graduação, pesquisadores, docentes e profissionais que desejem compreender melhor o papel do Linux na prática científica contemporânea.

A obra é construída de forma aberta e colaborativa. O uso de LaTeX permite manter uma estrutura editorial robusta, adequada à produção de material acadêmico. O uso de Git permite rastrear alterações, revisar contribuições, discutir modificações e preservar a história intelectual do projeto. Essa escolha metodológica reflete a própria filosofia do livro: aprender Linux não é apenas aprender comandos, mas aprender uma forma de organizar trabalho técnico de modo transparente, auditável e reprodutível.

O livro está organizado em partes progressivas. As primeiras seções apresentam fundamentos conceituais sobre sistemas operacionais, Linux, distribuições, instalação e formas de interação. Em seguida, o texto avança para uso prático, manipulação de arquivos, composição de comandos, automação, reprodutibilidade, integração com R e Python, infraestrutura computacional, HPC, containers, segurança, desempenho e filosofia Unix aplicada à ciência de dados.

% ============================================================
% Índice
% ============================================================

\cleardoublepage
\tableofcontents

% ============================================================
% Parte I
% ============================================================

\begin{refsection}
\part{Fundamentos}

\chapter{O que é um sistema operacional}

\section{Definição e função de um sistema operacional}

\topicos{
    \item Sistema operacional como mediador entre hardware e software.
    \item Abstração de recursos físicos.
    \item Interface para usuários e programas.
    \item Sistema operacional como camada de controle e coordenação.
}

\placeholder{Esta seção deve apresentar uma definição conceitual de sistema operacional, mostrando que ele organiza recursos computacionais e fornece abstrações para que programas possam executar sem manipular diretamente o hardware.}

\section{Gerenciamento de processos}

\topicos{
    \item O que é um processo.
    \item Diferença entre programa e processo.
    \item Escalonamento.
    \item Execução concorrente.
}

\placeholder{Esta seção deve explicar como programas em execução são representados pelo sistema operacional e como múltiplas tarefas podem compartilhar o processador.}

\section{Gerenciamento de memória}

\topicos{
    \item Memória física e memória virtual.
    \item Espaço de endereçamento.
    \item Paginação.
    \item Isolamento entre processos.
}

\placeholder{Esta seção deve introduzir a ideia de que cada processo opera em um espaço de memória abstraído pelo sistema operacional, o que permite isolamento, proteção e melhor uso dos recursos.}

\section{Sistema de arquivos}

\topicos{
    \item Arquivos como abstrações fundamentais.
    \item Diretórios.
    \item Permissões.
    \item Organização persistente de dados.
}

\placeholder{Esta seção deve mostrar que o sistema de arquivos é uma abstração central para armazenar, localizar e manipular dados, programas e configurações.}

\chapter{Como sistemas operacionais funcionam}

\section{Kernel e espaço de usuário}

\topicos{
    \item Kernel.
    \item User space.
    \item Separação de privilégios.
    \item Proteção do sistema.
}

\placeholder{Esta seção deve explicar a separação entre o núcleo do sistema operacional e os programas executados pelos usuários, mostrando por que essa separação é importante para estabilidade e segurança.}

\section{Chamadas de sistema}

\topicos{
    \item O que são system calls.
    \item Entrada e saída.
    \item Criação de processos.
    \item Acesso a arquivos e dispositivos.
}

\placeholder{Esta seção deve mostrar que programas comuns não acessam diretamente o hardware, mas solicitam serviços ao kernel por meio de chamadas de sistema.}

\section{Escalonamento e multitarefa}

\topicos{
    \item Compartilhamento da CPU.
    \item Estados de um processo.
    \item Prioridades.
    \item Impacto sobre desempenho.
}

\placeholder{Esta seção deve discutir como o sistema operacional decide qual processo executa em cada momento e como isso afeta programas científicos, scripts e pipelines de dados.}

\section{Entrada, saída e buffering}

\topicos{
    \item I/O de disco.
    \item I/O de terminal.
    \item Buffers.
    \item Gargalos de entrada e saída.
}

\placeholder{Esta seção deve conectar mecanismos internos de entrada e saída com problemas práticos de desempenho em análise de dados, especialmente quando arquivos grandes são manipulados.}

\chapter{Breve histórico do Linux e da filosofia Unix}

\section{Unix e sua influência}

\topicos{
    \item Origem do Unix.
    \item Sistemas multiusuário.
    \item Ferramentas pequenas e composáveis.
    \item Cultura de programação sistêmica.
}

\placeholder{Esta seção deve apresentar o Unix como matriz histórica e conceitual do Linux, destacando seu impacto na computação científica e na cultura de ferramentas composáveis.}

\section{O projeto GNU e o software livre}

\topicos{
    \item Origem do projeto GNU.
    \item Liberdades do software livre.
    \item Ferramentas GNU.
    \item Relação entre GNU e Linux.
}

\placeholder{Esta seção deve explicar a relevância do projeto GNU para o ecossistema Linux e introduzir a dimensão ética, técnica e política do software livre.}

\section{O kernel Linux}

\topicos{
    \item Surgimento do kernel Linux.
    \item Desenvolvimento colaborativo.
    \item Portabilidade.
    \item Adoção em servidores, pesquisa e dispositivos.
}

\placeholder{Esta seção deve apresentar o kernel Linux como projeto técnico e comunitário, mostrando sua importância em servidores, supercomputadores, sistemas embarcados e ambientes científicos.}

\section{Filosofia Unix}

\topicos{
    \item Fazer uma coisa e fazê-la bem.
    \item Composição de ferramentas.
    \item Texto como interface universal.
    \item Automação por scripts.
}

\placeholder{Esta seção deve introduzir a filosofia Unix como modelo de trabalho que permite construir soluções complexas por composição de ferramentas simples.}

\chapter{Arquitetura geral de um sistema Linux}

\section{Componentes principais}

\topicos{
    \item Kernel.
    \item Bibliotecas.
    \item Shell.
    \item Serviços.
    \item Aplicativos.
}

\placeholder{Esta seção deve oferecer uma visão em camadas de um sistema Linux, mostrando como diferentes componentes interagem para formar um ambiente funcional.}

\section{Sistema de inicialização}

\topicos{
    \item Boot.
    \item Bootloader.
    \item Init.
    \item systemd.
}

\placeholder{Esta seção deve explicar, em nível introdutório, como o sistema passa do firmware ao kernel e depois aos serviços de usuário.}

\section{Diretórios especiais}

\topicos{
    \item \texttt{/proc}.
    \item \texttt{/sys}.
    \item \texttt{/dev}.
    \item \texttt{/etc}.
}

\placeholder{Esta seção deve mostrar que alguns diretórios representam interfaces dinâmicas com o kernel, dispositivos e configuração do sistema.}

\chapter{Formas de interação com o sistema}

\section{Interface gráfica}

\topicos{
    \item Ambientes gráficos.
    \item Gerenciadores de janelas.
    \item Aplicativos gráficos.
    \item Vantagens e limitações da GUI.
}

\placeholder{Esta seção deve apresentar a interface gráfica como uma camada de interação amigável, útil para muitas tarefas, mas limitada para automação e composição.}

\section{Terminal e linha de comando}

\topicos{
    \item Terminal.
    \item Shell.
    \item Comandos.
    \item Scripts.
}

\placeholder{Esta seção deve introduzir o terminal como ambiente de interação textual, destacando sua importância para ciência de dados, automação e trabalho remoto.}

\section{Comparação entre CLI e GUI}

\topicos{
    \item Interatividade.
    \item Automação.
    \item Reprodutibilidade.
    \item Curva de aprendizado.
}

\placeholder{Esta seção deve comparar os paradigmas de interface e mostrar que CLI e GUI não são rivais, mas ferramentas complementares.}

\chapter{O que são distribuições Linux e como escolher}

\section{Definição de distribuição Linux}

\topicos{
    \item Kernel Linux.
    \item Userland.
    \item Gerenciador de pacotes.
    \item Políticas de atualização.
}

\placeholder{Esta seção deve explicar que uma distribuição Linux é um sistema operacional completo construído ao redor do kernel Linux.}

\section{Famílias de distribuições}

\topicos{
    \item Debian e derivadas.
    \item Red Hat e derivadas.
    \item Arch e derivadas.
    \item Distribuições científicas ou especializadas.
}

\placeholder{Esta seção deve apresentar famílias de distribuições e mostrar diferenças relevantes para estudantes, pesquisadores e usuários iniciantes.}

\section{Critérios de escolha}

\topicos{
    \item Estabilidade.
    \item Facilidade de uso.
    \item Comunidade.
    \item Documentação.
    \item Disponibilidade de pacotes esatísticos e científicos.
}

\placeholder{Esta seção deve orientar o estudante a escolher uma distribuição considerando objetivos, hardware, experiência prévia e suporte comunitário.}

\chapter{Processo geral de instalação do Linux}

\section{Preparação}

\topicos{
    \item Backup.
    \item Escolha da distribuição.
    \item Verificação da imagem ISO.
    \item Criação de mídia de instalação.
}

\placeholder{Esta seção deve descrever a preparação conceitual para instalação, enfatizando backup, segurança e verificação de integridade.}

\section{Boot e firmware}

\topicos{
    \item BIOS.
    \item UEFI.
    \item Ordem de boot.
    \item Secure Boot.
}

\placeholder{Esta seção deve explicar os conceitos gerais envolvidos no início da instalação, sem depender de uma distribuição específica.}

\section{Particionamento}

\topicos{
    \item Disco.
    \item Partição.
    \item Sistema de arquivos.
    \item Dual boot.
    \item Partição EFI.
}

\placeholder{Esta seção deve explicar decisões gerais de particionamento e seus riscos, preparando o estudante para compreender instaladores diferentes.}

\section{Instalação e pós-instalação}

\topicos{
    
    \item Instalação do bootloader.
    \item Particionamento
    \item Criação de usuário.
    \item Atualização do sistema.
    \item Instalação de pacotes essenciais.
}

\placeholder{Esta seção deve apresentar a lógica geral da instalação e as primeiras tarefas após reiniciar no sistema instalado.}

\chapter{Estrutura de diretórios e sistema de arquivos}

\section{Hierarquia padrão}

\topicos{
    \item \texttt{/home}.
    \item \texttt{/usr}.
    \item \texttt{/var}.
    \item \texttt{/tmp}.
    \item \texttt{/etc}.
}

\placeholder{Esta seção deve explicar a organização típica de diretórios em sistemas Linux e sua relevância para localizar dados, programas e configurações.}

\section{Permissões e propriedade}

\topicos{
    \item Usuário.
    \item Grupo.
    \item Permissões de leitura, escrita e execução.
    \item Comando \texttt{chmod}.
    \item Comando \texttt{chown}.
}

\placeholder{Esta seção deve introduzir o modelo de permissões como mecanismo de segurança e organização multiusuário.}

\section{Links e pontos de montagem}

\topicos{
    \item Links simbólicos.
    \item Hard links.
    \item Montagem.
    \item Dispositivos externos.
}

\placeholder{Esta seção deve explicar como arquivos e dispositivos podem ser referenciados e integrados à árvore de diretórios.}

\bibliografiaParte{recomendadosParteI}
\end{refsection}

% ============================================================
% Parte II
% ============================================================

\begin{refsection}
\part{Uso prático}

\chapter{Navegação e manipulação básica de arquivos}

\section{Navegação no terminal}

\topicos{
    \item \texttt{pwd}.
    \item \texttt{ls}.
    \item \texttt{cd}.
    \item Caminhos absolutos e relativos.
}

\placeholder{Esta seção deve ensinar a movimentação básica pela árvore de diretórios e desenvolver intuição espacial sobre o sistema de arquivos.}

\section{Criação, cópia, movimentação e remoção}

\topicos{
    \item \texttt{touch}.
    \item \texttt{mkdir}.
    \item \texttt{cp}.
    \item \texttt{mv}.
    \item \texttt{rm}.
}

\placeholder{Esta seção deve introduzir operações fundamentais de manipulação de arquivos, sempre destacando riscos, boas práticas e exemplos com diretórios de projeto.}

\section{Globbing e padrões}

\topicos{
    \item \texttt{*}.
    \item \texttt{?}.
    \item Listas de caracteres.
    \item Expansão pelo shell.
}

\placeholder{Esta seção deve explicar como o shell expande padrões de nomes de arquivos e como isso facilita operações em lote.}

\chapter{Visualização e inspeção de dados}

\section{Inspeção de arquivos}

\topicos{
    \item \texttt{cat}.
    \item \texttt{less}.
    \item \texttt{head}.
    \item \texttt{tail}.
}

\placeholder{Esta seção deve mostrar como examinar arquivos de texto, CSVs, logs e saídas de programas sem abrir editores pesados.}

\section{Arquivos grandes}

\topicos{
    \item Leitura parcial.
    \item Paginação.
    \item Evitar carregar tudo na memória.
    \item Estratégias de amostragem rápida.
}

\placeholder{Esta seção deve mostrar que arquivos grandes exigem estratégias diferentes de inspeção, particularmente em ciência de dados.}

\section{Encoding e caracteres}

\topicos{
    \item UTF-8.
    \item Problemas de acentuação.
    \item Arquivos vindos de diferentes sistemas.
    \item Conversão de codificação.
}

\placeholder{Esta seção deve introduzir problemas comuns de codificação textual que afetam bases de dados, planilhas exportadas e documentos científicos.}

\chapter{Pipes e composição de comandos}

\section{Entrada padrão, saída padrão e erro padrão}

\topicos{
    \item \texttt{stdin}.
    \item \texttt{stdout}.
    \item \texttt{stderr}.
    \item Redirecionamento.
}

\placeholder{Esta seção deve explicar o modelo de fluxo textual que permite conectar programas e registrar resultados.}

\section{Pipes}

\topicos{
    \item Operador \texttt{|}.
    \item Encadeamento de comandos.
    \item Composição funcional.
    \item Pipelines de dados.
}

\placeholder{Esta seção deve apresentar pipes como mecanismo central da filosofia Unix e como analogia operacional para pipelines de ciência de dados.}

\section{Redirecionamento}

\topicos{
    \item \texttt{>}.
    \item \texttt{>>}.
    \item \texttt{2>}.
    \item Arquivos de log.
}

\placeholder{Esta seção deve mostrar como salvar resultados, acumular saídas e separar erros de resultados úteis.}

\chapter{Filtragem e transformação de dados em escala}

\section{Filtragem textual}

\topicos{
    \item \texttt{grep}.
    \item Busca por padrões.
    \item Sensibilidade a maiúsculas e minúsculas.
    \item Contagem de ocorrências.
}

\placeholder{Esta seção deve introduzir buscas textuais em arquivos grandes e conectar o uso de filtros à exploração inicial de dados.}

\section{Seleção de colunas e campos}

\topicos{
    \item \texttt{cut}.
    \item Delimitadores.
    \item Campos.
    \item Arquivos tabulares.
}

\placeholder{Esta seção deve mostrar como extrair partes específicas de arquivos tabulares diretamente pelo terminal.}

\section{Ordenação, contagem e sumarização}

\topicos{
    \item \texttt{sort}.
    \item \texttt{uniq}.
    \item \texttt{wc}.
    \item Frequências simples.
}

\placeholder{Esta seção deve apresentar operações básicas de sumarização que podem ser executadas antes de abrir R ou Python.}

\section{Transformações com \texttt{awk} e \texttt{sed}}

\topicos{
    \item Campos e registros.
    \item Substituição textual.
    \item Transformações linha a linha.
    \item Pequenos programas de transformação.
}

\placeholder{Esta seção deve introduzir \texttt{awk} e \texttt{sed} como ferramentas para transformações rápidas, sem pretender substituir linguagens completas de análise.}

\chapter{Expressões regulares}

\section{Padrões e correspondência}

\topicos{
    \item Literais.
    \item Classes de caracteres.
    \item Quantificadores.
    \item Âncoras.
}

\placeholder{Esta seção deve apresentar expressões regulares como linguagem formal para descrever padrões textuais.}

\section{Captura e substituição}

\topicos{
    \item Grupos.
    \item Subexpressões.
    \item Substituição.
    \item Reorganização de texto.
}

\placeholder{Esta seção deve mostrar como expressões regulares permitem extrair, validar e reorganizar informações em arquivos textuais.}

\section{Aplicações em dados reais}

\topicos{
    \item Logs.
    \item Identificadores.
    \item Códigos de amostras.
    \item Dados semiestruturados.
}

\placeholder{Esta seção deve aplicar regex a problemas próximos da realidade de estudantes e pesquisadores.}

\chapter{One-liners como pipelines analíticos}

\section{O que é um one-liner}

\topicos{
    \item Comando curto.
    \item Composição.
    \item Operação reproduzível.
    \item Documentação por exemplos.
}

\placeholder{Esta seção deve definir one-liners como pequenas unidades de conhecimento prático, úteis para aprendizagem, automação e resolução rápida de problemas.}

\section{Legibilidade e manutenção}

\topicos{
    \item Quando usar one-liners.
    \item Quando transformar em script.
    \item Comentários.
    \item Clareza vs densidade.
}

\placeholder{Esta seção deve discutir os limites pedagógicos e práticos dos one-liners, evitando transformar concisão em obscuridade.}

\section{Biblioteca de one-liners}

\topicos{
    \item Organização por tema.
    \item Exemplos testáveis.
    \item Metadados.
    \item Contribuição colaborativa.
}

\placeholder{Esta seção deve conectar o livro a uma possível biblioteca colaborativa de comandos úteis para estudantes.}

\bibliografiaParte{recomendadosParteII}
\end{refsection}

% ============================================================
% Parte III
% ============================================================

\begin{refsection}
\part{Automação e reprodutibilidade}

\chapter{Scripts em shell}

\section{Por que escrever scripts}

\topicos{
    \item Repetição.
    \item Automação.
    \item Documentação executável.
    \item Redução de erro manual.
}

\placeholder{Esta seção deve mostrar que scripts transformam comandos interativos em procedimentos reprodutíveis.}

\section{Variáveis, argumentos e controle de fluxo}

\topicos{
    \item Variáveis.
    \item Argumentos posicionais.
    \item Condicionais.
    \item Laços.
}

\placeholder{Esta seção deve introduzir os elementos mínimos de programação em shell necessários para automatizar tarefas.}

\section{Boas práticas em shell scripts}

\topicos{
    \item Cabeçalho.
    \item Mensagens de erro.
    \item Verificação de argumentos.
    \item Organização.
}

\placeholder{Esta seção deve ensinar práticas defensivas para evitar scripts frágeis e difíceis de depurar.}

\chapter{Controle de versão e rastreabilidade}

\section{Introdução ao Git}

\topicos{
    \item Repositórios.
    \item Commits.
    \item Histórico.
    \item Branches.
}

\placeholder{Esta seção deve apresentar Git como ferramenta de rastreabilidade intelectual e técnica.}

\section{Git em projetos científicos}

\topicos{
    \item Versionamento de scripts.
    \item Versionamento de manuscritos.
    \item Registro de decisões.
    \item Colaboração.
}

\placeholder{Esta seção deve mostrar como Git se integra à escrita científica, à análise de dados e à colaboração em equipe.}

\section{Boas práticas de commits}

\topicos{
    \item Mensagens claras.
    \item Commits pequenos.
    \item Revisão por pares.
    \item Pull requests.
}

\placeholder{Esta seção deve conectar práticas de desenvolvimento de software a projetos acadêmicos e colaborativos.}

\chapter{Reprodutibilidade computacional}

\section{O que é reprodutibilidade}

\topicos{
    \item Reprodutibilidade.
    \item Replicabilidade.
    \item Transparência.
    \item Auditoria.
}

\placeholder{Esta seção deve definir reprodutibilidade computacional e explicar sua importância para ciência de dados e pesquisa científica.}

\section{Estrutura de projetos reproduzíveis}

\topicos{
    \item Diretórios padronizados.
    \item Dados brutos.
    \item Dados processados.
    \item Scripts.
    \item Resultados.
}

\placeholder{Esta seção deve propor uma organização mínima para projetos analíticos que possam ser compreendidos e reexecutados.}

\section{Logs, seeds e ambientes}

\topicos{
    \item Registro de execução.
    \item Aleatoriedade controlada.
    \item Versões de pacotes.
    \item Ambientes computacionais.
}

\placeholder{Esta seção deve apresentar elementos que afetam a capacidade de reproduzir resultados computacionais.}

\chapter{Automação de tarefas}

\section{Agendamento com cron}

\topicos{
    \item Tarefas periódicas.
    \item Sintaxe básica.
    \item Logs.
    \item Cuidados.
}

\placeholder{Esta seção deve introduzir execução periódica de comandos e scripts para tarefas repetitivas.}

\section{Automação local}

\topicos{
    \item Backups.
    \item Atualização de dados.
    \item Relatórios.
    \item Monitoramento.
}

\placeholder{Esta seção deve mostrar exemplos práticos de automação úteis para estudantes.}

\chapter{Build systems e workflows}

\section{Make}

\topicos{
    \item Alvos.
    \item Dependências.
    \item Regras.
    \item Reconstrução incremental.
}

\placeholder{Esta seção deve apresentar \texttt{make} como uma ferramenta geral para organizar dependências entre arquivos e comandos.}

\section{Workflows como grafos}

\topicos{
    \item DAGs.
    \item Etapas de análise.
    \item Dependências.
    \item Reexecução seletiva.
}

\placeholder{Esta seção deve conectar workflows computacionais à ideia de grafos acíclicos direcionados.}

\section{Snakemake e Nextflow}

\topicos{
    \item Motivação.
    \item Escalabilidade.
    \item Ambientes científicos.
    \item Bioinformática e ciência de dados.
}

\placeholder{Esta seção deve apresentar ferramentas modernas de workflow como evolução natural de scripts e Makefiles em projetos complexos.}

\bibliografiaParte{recomendadosParteIII}
\end{refsection}

% ============================================================
% Parte IV
% ============================================================

\begin{refsection}
\part{Integração com ciência de dados}

\chapter{Integração com R e Python}

\section{Execução de scripts}

\topicos{
    \item \texttt{Rscript}.
    \item \texttt{python}.
    \item Argumentos de linha de comando.
    \item Entrada e saída.
}

\placeholder{Esta seção deve mostrar como executar análises em R e Python a partir do terminal.}

\section{Pipelines híbridos}

\topicos{
    \item Shell como cola.
    \item Pré-processamento.
    \item Análise estatística.
    \item Relatórios.
}

\placeholder{Esta seção deve mostrar como Linux, R e Python podem funcionar como partes complementares de um mesmo pipeline.}

\section{Organização de projetos analíticos}

\topicos{
    \item Diretórios.
    \item Scripts.
    \item Dados.
    \item Resultados.
    \item Documentação.
}

\placeholder{Esta seção deve propor padrões de organização para projetos de análise de dados em ambiente Linux.}

\chapter{Gerenciamento de ambientes}

\section{Por que isolar ambientes}

\topicos{
    \item Dependências.
    \item Versões.
    \item Conflitos.
    \item Reprodutibilidade.
}

\placeholder{Esta seção deve explicar por que ambientes isolados são necessários em projetos com múltiplas bibliotecas e versões.}

\section{Ambientes Python}

\topicos{
    \item \texttt{venv}.
    \item \texttt{pip}.
    \item \texttt{requirements.txt}.
    \item Ambientes por projeto.
}

\placeholder{Esta seção deve introduzir práticas básicas de isolamento de ambientes Python.}

\section{Conda}

\topicos{
    \item Ambientes Conda.
    \item Pacotes científicos.
    \item Arquivos \texttt{environment.yml}.
    \item Vantagens e limitações.
}

\placeholder{Esta seção deve apresentar Conda como solução frequente em ciência de dados, bioinformática e computação científica.}

\chapter{Manipulação de dados grandes em ambiente local}

\section{Processamento por streaming}

\topicos{
    \item Leitura linha a linha.
    \item Evitar carregar tudo em memória.
    \item Composição de filtros.
    \item Pré-processamento no terminal.
}

\placeholder{Esta seção deve explicar por que operações baseadas em streaming são úteis para arquivos maiores que a memória disponível.}

\section{Amostragem e inspeção}

\topicos{
    \item Amostras aleatórias.
    \item Primeiras linhas.
    \item Últimas linhas.
    \item Estatísticas rápidas.
}

\placeholder{Esta seção deve ensinar estratégias de exploração preliminar de arquivos grandes.}

\section{Compressão}

\topicos{
    \item \texttt{gzip}.
    \item \texttt{xz}.
    \item \texttt{zcat}.
    \item Pipelines com arquivos comprimidos.
}

\placeholder{Esta seção deve mostrar como trabalhar com dados comprimidos sem necessariamente descompactar tudo em disco.}

\chapter{Introdução a bancos de dados via terminal}

\section{SQLite}

\topicos{
    \item Banco de dados em arquivo único.
    \item Tabelas.
    \item Consultas SQL.
    \item Uso local.
}

\placeholder{Esta seção deve apresentar SQLite como uma ponte simples entre arquivos tabulares e bancos de dados relacionais.}

\section{Consultas básicas}

\topicos{
    \item \texttt{SELECT}.
    \item \texttt{WHERE}.
    \item \texttt{GROUP BY}.
    \item \texttt{ORDER BY}.
}

\placeholder{Esta seção deve introduzir consultas SQL essenciais para estudantes de dados.}

\section{Integração com pipelines}

\topicos{
    \item Importação de CSV.
    \item Exportação de resultados.
    \item Uso com shell.
    \item Uso com R e Python.
}

\placeholder{Esta seção deve mostrar como bancos de dados locais podem participar de fluxos reprodutíveis de análise.}

\bibliografiaParte{recomendadosParteIV}
\end{refsection}

% ============================================================
% Parte V
% ============================================================

\begin{refsection}
\part{Sistemas e infraestrutura}

\chapter{Processos e monitoramento}

\section{Processos em execução}

\topicos{
    \item \texttt{ps}.
    \item PID.
    \item Usuário.
    \item Estado do processo.
}

\placeholder{Esta seção deve explicar como observar programas em execução e identificar processos relevantes.}

\section{Monitoramento interativo}

\topicos{
    \item \texttt{top}.
    \item \texttt{htop}.
    \item CPU.
    \item Memória.
}

\placeholder{Esta seção deve mostrar como acompanhar uso de recursos durante análises e scripts longos.}

\section{Encerramento de processos}

\topicos{
    \item Sinais.
    \item \texttt{kill}.
    \item \texttt{SIGTERM}.
    \item \texttt{SIGKILL}.
}

\placeholder{Esta seção deve explicar como encerrar processos de modo seguro e quando medidas mais fortes são necessárias.}

\chapter{Gerenciamento de pacotes}

\section{Conceitos fundamentais}

\topicos{
    \item Pacotes.
    \item Repositórios.
    \item Dependências.
    \item Atualizações.
}

\placeholder{Esta seção deve explicar o papel do gerenciador de pacotes na instalação e manutenção de software.}

\section{Gerenciadores comuns}

\topicos{
    \item \texttt{apt}.
    \item \texttt{dnf}.
    \item \texttt{pacman}.
    \item Diferenças entre distribuições.
}

\placeholder{Esta seção deve apresentar gerenciadores de pacotes de diferentes famílias de distribuições.}

\section{Binários e compilação}

\topicos{
    \item Pacotes binários.
    \item Código-fonte.
    \item Compilação.
    \item Dependências de desenvolvimento.
}

\placeholder{Esta seção deve explicar quando é necessário compilar software e quais riscos isso introduz.}

\chapter{Redes e acesso remoto}

\section{Conceitos básicos de rede}

\topicos{
    \item IP.
    \item Portas.
    \item Protocolos.
    \item Cliente e servidor.
}

\placeholder{Esta seção deve introduzir conceitos mínimos de rede necessários para uso remoto de sistemas Linux.}

\section{SSH}

\topicos{
    \item Login remoto.
    \item Chaves SSH.
    \item Segurança.
    \item Execução remota de comandos.
}

\placeholder{Esta seção deve apresentar SSH como ferramenta essencial para servidores, clusters e colaboração técnica.}

\section{Transferência de arquivos}

\topicos{
    \item \texttt{scp}.
    \item \texttt{rsync}.
    \item Sincronização.
    \item Transferências incrementais.
}

\placeholder{Esta seção deve mostrar métodos seguros e eficientes de transferir dados entre máquinas.}

\chapter{Computação em clusters e HPC}

\section{Do computador pessoal ao cluster}

\topicos{
    \item Recursos compartilhados.
    \item Nós computacionais.
    \item Filas.
    \item Jobs.
}

\placeholder{Esta seção deve explicar a diferença entre executar análises localmente e usar infraestrutura compartilhada de alto desempenho.}

\section{SLURM}

\topicos{
    \item Submissão de jobs.
    \item Scripts batch.
    \item Alocação de recursos.
    \item Monitoramento.
}

\placeholder{Esta seção deve introduzir SLURM como sistema comum de gerenciamento de jobs em clusters.}

\section{Boas práticas em HPC}

\topicos{
    \item Solicitação adequada de recursos.
    \item Logs.
    \item Organização de resultados.
    \item Respeito ao ambiente compartilhado.
}

\placeholder{Esta seção deve discutir responsabilidade e eficiência no uso de clusters.}

\chapter{Containers e ambientes reprodutíveis}

\section{O que são containers}

\topicos{
    \item Isolamento.
    \item Imagens.
    \item Execução portável.
    \item Diferença entre container e máquina virtual.
}

\placeholder{Esta seção deve explicar containers como mecanismo de empacotamento de ambientes computacionais.}

\section{Docker}

\topicos{
    \item Imagens.
    \item Containers.
    \item Dockerfile.
    \item Volumes.
}

\placeholder{Esta seção deve introduzir conceitos básicos de Docker e sua utilidade para ciência de dados.}

\section{Containers em ciência}

\topicos{
    \item Reprodutibilidade.
    \item Ambientes complexos.
    \item Distribuição de pipelines.
    \item Limitações e segurança.
}

\placeholder{Esta seção deve discutir containers como ferramenta de infraestrutura científica, sem esconder seus limites.}

\bibliografiaParte{recomendadosParteV}
\end{refsection}

% ============================================================
% Parte VI
% ============================================================

\begin{refsection}
\part{Prática avançada}

\chapter{Debugging e análise de erros}

\section{Como ler erros}

\topicos{
    \item Mensagens de erro.
    \item Código de saída.
    \item Logs.
    \item Diagnóstico incremental.
}

\placeholder{Esta seção deve ensinar o estudante a tratar erros como fontes de informação, não como eventos misteriosos.}

\section{Saída de erro e logs}

\topicos{
    \item \texttt{stderr}.
    \item Redirecionamento.
    \item Arquivos de log.
    \item Rastreamento de execução.
}

\placeholder{Esta seção deve aprofundar a distinção entre resultados e mensagens diagnósticas.}

\section{Estratégias de depuração}

\topicos{
    \item Reduzir o problema.
    \item Testar partes isoladas.
    \item Verificar entradas.
    \item Registrar passos intermediários.
}

\placeholder{Esta seção deve apresentar um método geral de depuração aplicável a scripts, comandos e pipelines.}

\chapter{Performance e otimização}

\section{Gargalos comuns}

\topicos{
    \item CPU.
    \item Memória.
    \item Disco.
    \item Rede.
}

\placeholder{Esta seção deve ensinar a distinguir gargalos computacionais e evitar otimizações prematuras.}

\section{Paralelismo simples}

\topicos{
    \item \texttt{xargs -P}.
    \item GNU Parallel.
    \item Tarefas independentes.
    \item Riscos de concorrência.
}

\placeholder{Esta seção deve introduzir paralelismo prático para tarefas independentes comuns em ciência de dados.}

\section{Profiling básico}

\topicos{
    \item Tempo de execução.
    \item Uso de memória.
    \item Comparação de abordagens.
    \item Medição antes de otimizar.
}

\placeholder{Esta seção deve mostrar que decisões de otimização devem ser guiadas por medidas.}

\chapter{Segurança básica}

\section{Usuários, grupos e privilégios}

\topicos{
    \item Usuários.
    \item Grupos.
    \item Permissões.
    \item \texttt{sudo}.
}

\placeholder{Esta seção deve apresentar o modelo básico de segurança de um sistema Linux multiusuário.}

\section{Boas práticas}

\topicos{
    \item Atualizações.
    \item Senhas.
    \item Chaves SSH.
    \item Cuidado com comandos destrutivos.
}

\placeholder{Esta seção deve ensinar práticas mínimas de segurança para estudantes que começam a usar Linux.}

\section{Segurança em dados científicos}

\topicos{
    \item Dados sensíveis.
    \item Controle de acesso.
    \item Backup.
    \item Compartilhamento responsável.
}

\placeholder{Esta seção deve conectar segurança computacional à responsabilidade científica e ética no tratamento de dados.}

\bibliografiaParte{recomendadosParteVI}
\end{refsection}

\begin{refsection}
\part{Filosofia UNIX}

\chapter{Filosofia Unix aplicada à ciência de dados}

\section{Composição}

\topicos{
    \item Ferramentas pequenas.
    \item Interfaces textuais.
    \item Pipelines.
    \item Modularidade.
}

\placeholder{Esta seção deve sintetizar a ideia de que tarefas complexas podem ser decompostas em operações simples conectadas.}

\section{Minimalismo operacional}

\topicos{
    \item Soluções pequenas.
    \item Redução de dependências.
    \item Clareza.
    \item Reutilização.
}

\placeholder{Esta seção deve discutir como o minimalismo técnico pode produzir fluxos de trabalho mais robustos e compreensíveis.}

\section{Pipelines como grafos}

\topicos{
    \item Nós.
    \item Arestas.
    \item Transformações.
    \item Fluxos de dados.
}

\placeholder{Esta seção deve reinterpretar pipelines Unix como grafos computacionais, aproximando Linux de conceitos modernos de workflows de dados.}

\chapter{Do laptop ao ambiente de produção}

\section{Escala de maturidade computacional}

\topicos{
    \item Uso individual.
    \item Projetos reproduzíveis.
    \item Servidores.
    \item Clusters.
    \item Produção.
}

\placeholder{Esta seção deve mostrar como as práticas aprendidas no livro podem evoluir para ambientes mais profissionais.}

\section{Organização de projetos de longo prazo}

\topicos{
    \item Estrutura.
    \item Documentação.
    \item Automação.
    \item Manutenção.
}

\placeholder{Esta seção deve discutir como manter projetos científicos e computacionais compreensíveis ao longo do tempo.}

\section{Boas práticas profissionais}

\topicos{
    \item Revisão.
    \item Testes.
    \item Integração contínua.
    \item Distribuição.
}

\placeholder{Esta seção deve apresentar práticas profissionais que podem ser incorporadas gradualmente por estudantes.}

\bibliografiaParte{recomendadosParteVII}
\end{refsection}

% ============================================================
% Próximos passos
% ============================================================

\begin{refsection}
\chapter{Próximos passos}

\section{Como continuar aprendendo Linux}

\topicos{
    \item Prática cotidiana.
    \item Leitura de manuais.
    \item Exploração de projetos reais.
    \item Participação em comunidades.
}

\placeholder{Esta seção deve orientar o estudante a transformar Linux em ferramenta cotidiana de estudo e trabalho.}

\section{Projetos sugeridos}

\topicos{
    \item Organizar um projeto de análise reproduzível.
    \item Criar scripts de automação.
    \item Montar um ambiente científico em Linux.
    \item Contribuir com software livre.
}

\placeholder{Esta seção deve propor atividades práticas que consolidem os conhecimentos do livro.}

\section{Trilhas de aprofundamento}

\topicos{
    \item Administração de sistemas.
    \item Ciência de dados reprodutível.
    \item HPC.
    \item Engenharia de dados.
    \item Desenvolvimento de software.
}

\placeholder{Esta seção deve apresentar caminhos possíveis após a leitura do livro, conectando Linux a trajetórias acadêmicas e profissionais.}

\section{Participação no GULECD/UFPR}

\topicos{
    \item Grupos de estudo.
    \item Oficinas.
    \item Projetos colaborativos.
    \item Apoio a novos estudantes.
}

\placeholder{Esta seção deve convidar o estudante a participar da comunidade e transformar aprendizado individual em prática coletiva.}

\cleardoublepage
\section*{Bibliografia do Capítulo}
\addcontentsline{toc}{section}{Bibliografia do Capítulo}
\printbibliography[
    heading=none,
    notcategory=recomendadosProximosPassos
]

\section*{Leitura Recomendada}
\addcontentsline{toc}{section}{Leitura Recomendada}
\nocite{lacroixLinuxMintEssentials2014,jrLinuxCommandLine2012}
\printbibliography[
    heading=none,
    category=recomendadosProximosPassos
]
\end{refsection}

% ============================================================
% Glossário
% ============================================================

\cleardoublepage
\printglossary[title={Glossário},toctitle={Glossário}]

% ============================================================
% Sobre o GULECD/UFPR
% ============================================================

\cleardoublepage
\chapter*{Sobre o Grupo de Usuários Linux e Software Livre do Curso de Estatística e Ciência de Dados da UFPR}
\addcontentsline{toc}{chapter}{Sobre o GULECD/UFPR}

O Grupo de Usuários Linux e Software Livre do Curso de Estatística e Ciência de Dados da UFPR, também referido como GULECD/UFPR, é uma iniciativa estudantil, colaborativa e acadêmica voltada à promoção do uso de Linux, software livre, ciência aberta, computação científica e boas práticas computacionais no contexto da formação em Estatística e Ciência de Dados.

O grupo nasce da percepção de que a formação contemporânea em estatística, ciência de dados e pesquisa quantitativa exige mais do que domínio de métodos matemáticos e estatísticos. Exige também capacidade de operar ambientes computacionais reais, compreender sistemas, automatizar tarefas, trabalhar de forma reprodutível, colaborar em projetos versionados e usar ferramentas abertas de modo crítico e competente.

Entre suas atividades potenciais estão oficinas, cursos, grupos de estudo, produção de material didático, apoio a estudantes iniciantes, desenvolvimento de pequenos projetos de software, promoção de boas práticas de reprodutibilidade e articulação com docentes, técnicos e demais membros da comunidade universitária.

Este livro é uma das expressões desse projeto formativo. Ele busca oferecer aos estudantes um caminho de entrada ao Linux que seja tecnicamente consistente, pedagogicamente acessível e alinhado às necessidades reais de quem trabalha com dados, modelos, programação, pesquisa científica e colaboração aberta.

% ============================================================
% Fim do documento
% ============================================================

\end{document}